\documentclass[10pt,a4paper]{amsart}
\usepackage[margin=19mm]{geometry}
\usepackage{amsmath,amssymb,mathtools}
\usepackage[scr=rsfs,cal=euler]{mathalfa}
\usepackage{xfrac}
\usepackage[unicode,hidelinks,bookmarks=false]{hyperref}
\newcommand{\ie}{i.e.\ }
\newcommand{\cf}{cf.\ }
\newcommand{\resp}{respectively\ }
\newcommand{\Subsection}{\subsection}
\providecommand{\nobreakdash}{\nobreak\hbox{-}}
\setlength{\emergencystretch}{2em}
\multlinegap0pt
\usepackage{amssymb}
\usepackage{bm}
\usepackage[shortlabels]{enumitem}
\usepackage{xcolor}
\newtheorem{thm}{Theorem}
\def\R{\mathbb R}
\def\Z{\mathbb Z}
\title{Rigidity and non-rigidity of the stable norm on $T^n$}
\author{Fernando C. Marques, Andr\'e Neves, Ao Sun}
\date{Source: arXiv:2608.15376v1 (CC BY 4.0)}
\begin{document}
\maketitle

\noindent\emph{Source and licence.} Rigidity and non-rigidity of the stable norm on $T^n$. Version arXiv:2608.15376v1 (CC BY 4.0), distributed by arXiv under the Creative Commons Attribution 4.0 International licence, http://creativecommons.org/licenses/by/4.0/, as recorded in the arXiv licence metadata for this exact version and shown on the abstract page https://arxiv.org/abs/2608.15376v1. Full licence notice: \texttt{licenses/arxiv-2608.15376-CC-BY-4.0.txt}.\par\smallskip

\noindent\textit{Editorial scope.} Counted unit: the article's thm thm (source0.tex lines 1370--1384). The article's complete original proof of it is delivered with it (source0.tex lines 1386--1478, one \texttt{proof} environment of 93 lines). No mathematics is written, repaired or re-derived: the statement and the proof are the article's own lines, in the article's own notation, and no retained line is substituted or re-flowed.  No further source-local context is needed: every cross-reference of every retained line resolves inside the two delivered spans.  Nothing was omitted from any retained span, so the package carries no omission record.  No retained line contains a citation, so the wrapper carries no bibliography and no bibliography entry.  No retained line contains a figure environment, an image-inclusion command or drawing code, so no image is delivered and the package declares no asset dependency.  Editorial formatting only, in addition: an AMS \texttt{amsart} preamble that restates the article's own declarations and macros used by the retained lines, namely the environments \texttt{thm} on separate counters, together with the standard packages those lines need.  The retained environments print in this document as Theorem 1 in order of appearance, whereas the article prints the same items with the numbering of its own full document.  The complete native source of the article as distributed in the arXiv e-print for this exact version is bundled unchanged as \texttt{sources/207741/source0.tex}.
		\begin{thm}
			Let 
			\begin{equation}
				V(u)=-\frac{1}{2}(f(u))^2,
			\end{equation}
			where 
			\begin{enumerate}
				\item $f$ is $1$-periodic,
				\item $f\geq 0$,
				\item $f(0)=0$.
			\end{enumerate}
			(For example, we can choose $f(u)=\epsilon(\sin(\pi u))^2$ for any $\epsilon>0$.)
			
			Then for every $\alpha \in \mathbb R^n$ there exists a foliation of $T^{n+1}$ by graphs of minimal solutions of slope $\alpha$.
		\end{thm}

		\begin{proof}
			By a direct calculation, $-V'(u)=f'(u)f(u)$.
			
			{\bf Case 1. $\alpha=0$.} Define $\varphi_t(z)$ to be the flow in the universal cover $\R^n$ of the equation 
			\[
			\frac{d}{dt}\varphi_t(z)=f(z),\quad \varphi_0(z)=z.
			\]
			Note that $\varphi_{t}(z+1)=\varphi_t(z)+1$.
			Define a map $\Phi:T^{n+1}\to T^{n+1}$ by 
			\begin{equation}
				\Phi(x,z)=(x,\varphi_{x_1}(z) \text{ mod $\Z$}).   
			\end{equation}
			Moreover, $\Phi$ is a self-diffeomorphism because for a fixed $t$ $\varphi_t$ is a self-diffeomorphism of $\R$. Thus, for $z\in[0,1)$, the graphs of the functions
			\begin{equation}
				U_z(x):=\varphi_{x_1}(z)
			\end{equation}
			form a foliation of $T^{n+1}$. Direct calculation shows that
			\begin{equation}
				\begin{split}
					\Delta (U_z(x))
					&=\Delta(\varphi_{x_1}(z))
					=
					\frac{d^2}{dx_1^2}\varphi_{x_1}(z)
					=\frac{d}{dx_1}f(\varphi_{x_1}(z))
					=
					f'(\varphi_{x_1}(z))f(\varphi_{x_1}(z))
					\\
					&= f'(U_z(x))f(U_z(x)).
				\end{split}
			\end{equation}
			So $U_z$ is a family of functions that solve the equation $\Delta U_z =- V' (U_z)$. Moreover, because $f(u)$ vanishes at $u=0$ and $u=1$, each $U_z(x)=\varphi_{x_1}(z)$ is uniformly bounded, so they have slope $\alpha=0$.
			
			In addition, $U_z$ is not only a solution, but a global minimizer of the potential. In fact, if $F$ satisfies $F'=f$, then for arbitrary smooth function $v$, we have
			\begin{equation}
				\frac{1}{2}|\nabla v|^2+\frac{1}{2}f(v)^2=\frac{1}{2}|\nabla v-f(v)e_1|^2+\partial_{x_1}F(v).
			\end{equation}
			For the minimizer $U_z$ that we just constructed, $|\nabla U_z-f(U_z)e_1|^2=0$. Thus, for $\Omega$ a bounded domain in $\R^n$ in the universal cover, and for any $v=U_z$ on $\partial\Omega$, by the divergence theorem, we have
			\[
			\int_\Omega \partial_{x_1}F(v)dx=\int_\Omega \partial_{x_1}F(U_z)dx.
			\]
			This implies that
			\begin{equation}
				\begin{split}
					\int_{\Omega} \left(
					\frac{1}{2}|\nabla v|^2+\frac{1}{2}f(v)^2
					\right)dx
					-
					\int_{\Omega}& \left(
					\frac{1}{2}|\nabla U_z|^2+\frac{1}{2}f(U_z)^2
					\right)dx
					\\&=
					\int_{\Omega} \left(
					\frac{1}{2}|\nabla v-f(v)e_1|^2    \right)dx\geq 0.
				\end{split}
			\end{equation}
			So $U_z$ is globally minimizing.
			
			\medskip
			
			
			{\bf Case 2. $\alpha\neq 0$.} 
			We want to find an increasing function $h_\alpha$ and define 
			\begin{equation}
				U_{z}(x)=h_{\alpha}(\alpha\cdot x+z).
			\end{equation}
			Note that 
			\[
			\begin{split}
				\nabla U_z&=\alpha  h_\alpha',
				\\
				\Delta U_z&=|\alpha|^2 h_\alpha''.
			\end{split}
			\]
			So we want to solve $h_\alpha''=|\alpha|^{-2}f'(h_\alpha)f(h_\alpha)$. This implies that
			\[
			((h_\alpha')^2)'=(|\alpha|^{-2}f^2(h_\alpha))'.
			\]
			Hence, we need to solve the equation
			\begin{equation}
				h_{\alpha}'(t)=|\alpha|^{-1}\sqrt{f^2(h_\alpha(t))+C},\qquad h_\alpha(0)=0.
			\end{equation}
			We claim that we can choose $C>0$ so that $h_{\alpha}(t+1)=h_{\alpha}(t)+1$. In fact, when $C>|\alpha|^2$, $h'_\alpha(t)> 1$ and hence $h_{\alpha}(t+1)>h_{\alpha}(t)+1$; when $C=0$, the solution $h_{\alpha}(t)\equiv 0$, so $h_{\alpha}(t+1)<h_{\alpha}(t)+1$. By the continuity dependence of ODE, there must be some $C_\alpha>0$ that makes $h_{\alpha}(t+1)=h_{\alpha}(t)+1$. With this choice of $h_{\alpha}$, we have
			\begin{equation}
				\Delta U_z=-V'(U_z).
			\end{equation}
			Since $h_\alpha$ is increasing, the graph of $U_z$ forms a foliation. Also, since $h_{\alpha}(t)-t$ is $1$-periodic, the graph of $U_z$ has slope $\alpha$.
			
			Finally, we show that $U_z$ is a globally minimizing solution. Similar to the zero slope case, we define $H'(u)=\sqrt{f(u)^2+C_\alpha}$. Let $e=\alpha/|\alpha|$. For any smooth function $v$,
			\begin{equation}
				\frac{1}{2}|\nabla v|^2+\frac{1}{2}f(v)^2=\frac{1}{2}|\nabla v-e\sqrt{f(u)^2+C_\alpha}|^2+e\cdot \nabla H(v)-\frac{1}{2}C_\alpha.    
			\end{equation}
			Then the same argument as the zero slope case shows that $U_z$ is globally minimizing.
		\end{proof}

\end{document}
